% ============================================
% Chapter 5: Results and Analysis
% ============================================

\chapter{Results and Analysis}

This chapter presents the dynamic scheduling benchmark results from \texttt{experiments/dynamic\_benchmark/}. All algorithms under the same scale share an identical task-generation seed; small scale is averaged over two seeds, while medium and large scales use one seed each. All time-based metrics are in simulation seconds.

\section{Small-Scale Results}

\subsection{Performance Table}

The small-scale configuration uses 3 vehicles, 1 charging station, and a task budget of 20. Table~\ref{tab:small-results} reports the mean over seed~1 and seed~2 (11--12 tasks settled per run).

\begin{table}[H]
    \centering
    \caption{Small-Scale Dynamic Benchmark Results (mean of 2 seeds)}
    \label{tab:small-results}
    \reporttableset{0.98\textwidth}
    \begin{reporttabular}{>{\raggedright\arraybackslash}l cccc}
        \toprule
        \textbf{Algorithm} & \textbf{Comp.(\%)} & \textbf{On-time(\%)} & \textbf{Total Score} & \textbf{Avg Time (s)} \\
        \midrule
        RL Batch                 & 100.0 & 73.3 & $\mathbf{1262.4}$ & 5974 \\
        SA Score Search          & 100.0 & 72.5 & 1243.9 & 4190 \\
        DFS Score Search         & 100.0 & \textbf{81.7} & 1242.4 & 4888 \\
        Nearest Task First       & 100.0 & 68.3 & 1238.4 & 5177 \\
        Regional Planning        & 100.0 & 72.5 & 1171.0 & 5304 \\
        Priority First           & 100.0 & 72.5 & 1166.4 & 4564 \\
        Heaviest Task First      & 100.0 & 63.3 & 1066.1 & 5531 \\
        Cluster Auction MAS      & 100.0 & 64.2 & 1037.4 & 5155 \\
        Relay Handoff            & 100.0 & 64.2 & 1030.6 & 5793 \\
        Earliest Deadline First  & 100.0 & 72.5 & 940.7 & 4717 \\
        \bottomrule
    \end{reporttabular}
\end{table}

\begin{figure}[H]
    \centering
    \includegraphics[width=0.92\textwidth]{chapters/comparison0.png}
    \caption{Small-scale total score comparison seed~1.}
    \label{fig:small-score}
\end{figure}



\begin{figure}[H]
    \centering
    \includegraphics[width=0.92\textwidth]{chapters/comparison00.png}
    \caption{Small-scale total score comparison seed~2.}
    \label{fig:small-score}
\end{figure}




\subsection{Analysis}

\begin{itemize}[leftmargin=4em]
    \item All ten algorithms achieve \textbf{100\% completion} on both seeds; differentiation comes from total score and on-time rate.
    \item Search-based batch loaders (\textbf{RL Batch}, \textbf{SA Score Search}, \textbf{DFS Score Search}) occupy the top tier (1242--1262 mean score). \textbf{DFS Score Search} also yields the highest mean on-time rate (81.7\%).
    \item \textbf{Earliest Deadline First} scores lowest (940.7 mean), indicating that deadline-only sorting is insufficient when batch loading and battery constraints dominate.
    \item Seed~2 is more volatile: \textbf{Earliest Deadline First} drops to 699.1 while \textbf{DFS Score Search} reaches 1142.6 with 83.3\% on-time, confirming sensitivity to task sequence even under shared infrastructure.
    \item Cooperative strategies (\textbf{Relay Handoff}, \textbf{Cluster Auction MAS}) do not outperform greedy baselines at this scale, likely because the fleet is small and handoff overhead is not yet justified.
\end{itemize}

\section{Medium-Scale Results}

\subsection{Performance Table}

The medium-scale run uses 6 vehicles, 2 charging stations, and settles 49 tasks. Results from seed~1 are shown in Table~\ref{tab:medium-results}.

\begin{table}[H]
    \centering
    \caption{Medium-Scale Dynamic Benchmark Results (seed~1, 49 tasks)}
    \label{tab:medium-results}
    \reporttableset{0.98\textwidth}
    \begin{reporttabular}{>{\raggedright\arraybackslash}l cccc}
        \toprule
        \textbf{Algorithm} & \textbf{Comp.(\%)} & \textbf{On-time(\%)} & \textbf{Total Score} & \textbf{Avg Time (s)} \\
        \midrule
        DFS Score Search         & 100.0 & 53.1 & $\mathbf{3824.3}$ & 6990 \\
        RL Batch                 & 100.0 & 53.1 & 3795.3 & 6311 \\
        SA Score Search          & 100.0 & 51.0 & 3716.2 & 7353 \\
        Nearest Task First       & 100.0 & \textbf{57.1} & 3377.2 & 7118 \\
        Relay Handoff            & 100.0 & 51.0 & 2913.1 & 7503 \\
        Priority First           & 100.0 & 40.8 & 2824.1 & 8869 \\
        Regional Planning        & 100.0 & 49.0 & 2385.3 & 7099 \\
        Cluster Auction MAS      & 100.0 & 49.0 & 2385.3 & 7145 \\
        Heaviest Task First      & 100.0 & 40.8 & 2112.6 & 9240 \\
        Earliest Deadline First  & 100.0 & 46.9 & 1716.3 & 8190 \\
        \bottomrule
    \end{reporttabular}
\end{table}

\begin{figure}[H]
    \centering
    \includegraphics[width=0.92\textwidth]{chapters/comparison1.png}
    \caption{Medium-scale total score comparison (seed~1, 49 tasks).}
    \label{fig:medium-score}
\end{figure}

\subsection{Analysis}

\begin{itemize}[leftmargin=4em]
    \item Completion remains \textbf{100\%} for all algorithms, but on-time rates fall to 41--57\%, reflecting tighter deadlines relative to fleet capacity.
    \item \textbf{DFS Score Search} (3824.3) and \textbf{RL Batch} (3795.3) clearly outperform greedy baselines, validating explicit batch-score optimization under higher load.
    \item Simple greedy rules diverge: \textbf{Nearest Task First} stays competitive (3377.2) while \textbf{Earliest Deadline First} collapses to 1716.3 --- deadline sorting alone cannot compensate for poor batch composition.
    \item Multi-vehicle coordinators (\textbf{Regional Planning}, \textbf{Cluster Auction MAS}) score 2385.3 with identical outcomes on this seed, suggesting similar package-decomposition behavior but no advantage over search-based per-vehicle batching here.
    \item The score gap between best and worst widens to 2108 points (3824.3 vs.\ 1716.3), showing that algorithm choice matters substantially at medium scale.
\end{itemize}

\section{Large-Scale Results}

\subsection{Performance Table}

The large-scale run uses 12 vehicles, 4 charging stations, and settles 32 tasks within the generation window. Table~\ref{tab:large-results} lists seed~1 results.

\begin{table}[H]
    \centering
    \caption{Large-Scale Dynamic Benchmark Results (seed~1, 32 tasks)}
    \label{tab:large-results}
    \reporttableset{0.98\textwidth}
    \begin{reporttabular}{>{\raggedright\arraybackslash}l cccc}
        \toprule
        \textbf{Algorithm} & \textbf{Comp.(\%)} & \textbf{On-time(\%)} & \textbf{Total Score} & \textbf{Avg Time (s)} \\
        \midrule
        Cluster Auction MAS      & 96.9 & 96.8 & $\mathbf{3278.7}$ & 1194 \\
        SA Score Search          & 100.0 & \textbf{96.9} & 3221.3 & 1089 \\
        Regional Planning        & 100.0 & 96.9 & 3045.0 & 1530 \\
        Heaviest Task First      & 100.0 & 93.8 & 2982.8 & 1395 \\
        RL Batch                 & 100.0 & 93.8 & 2907.9 & 1305 \\
        Earliest Deadline First  & 100.0 & 93.8 & 2907.9 & 1323 \\
        DFS Score Search         & 100.0 & 93.8 & 2907.9 & 1305 \\
        Priority First           & 100.0 & 93.8 & 2727.1 & 1377 \\
        Nearest Task First       & 100.0 & 96.9 & 2652.6 & 1368 \\
        Relay Handoff            & 100.0 & 93.8 & 2244.0 & 1764 \\
        \bottomrule
    \end{reporttabular}
\end{table}

\begin{figure}[H]
    \centering
    \includegraphics[width=0.92\textwidth]{chapters/comparison2.png}
    \caption{Large-scale total score comparison (seed~1, 32 tasks).}
    \label{fig:large-score}
\end{figure}

\subsection{Analysis}

\begin{itemize}[leftmargin=4em]
    \item Most algorithms complete all 32 tasks; \textbf{Cluster Auction MAS} stops at 31/32 (96.9\% completion) with one timeout, yet still achieves the highest total score (3278.7) on this seed.
    \item \textbf{SA Score Search} balances score (3221.3) with full completion and the joint-best on-time rate (96.9\%).
    \item \textbf{Relay Handoff} scores lowest (2244.0) despite 100\% completion --- relay negotiation adds route overhead without sufficient cargo-rebalancing benefit on this instance.
    \item Several batch-search variants (\textbf{RL Batch}, \textbf{DFS Score Search}, \textbf{Earliest Deadline First}) converge to identical scores (2907.9), indicating diminishing returns of search when the pending pool is small at each decision step.
    \item Compared with medium scale, absolute scores drop because fewer tasks are settled before the generation cutoff, highlighting that scale configuration affects both difficulty and score magnitude.
\end{itemize}

\section{Comparison with Static Gurobi Solver}


\begin{figure}[H]
    \centering
    \includegraphics[width=0.92\textwidth]{chapters/gurobi.png}
    \caption{static result for gurobi.}
    \label{fig:small-score}
\end{figure}



As an supplementary baseline, we also ran the static omniscient solver (Gurobi) on the same task-generation seeds. Table~\ref{tab:static-gurobi} summarizes those results.

\begin{table}[H]
    \centering
    \caption{Static Omniscient Solver Results (Gurobi)}
    \label{tab:static-gurobi}
    \begin{reporttabular}{lccc}
        \toprule
        \textbf{Scale} & \textbf{Total Score} & \textbf{Completed / Total} & \textbf{Completion Rate} \\
        \midrule
        Small  & 877.9  & 10 / 10 & 100.0\% \\
        Medium & 3003.6 & 20 / 20 & 100.0\% \\
        Large  & 4407.0 & 38 / 39 & 97.4\%  \\
        \bottomrule
    \end{reporttabular}
\end{table}

The static solver results differ substantially from the dynamic benchmark and do not consistently dominate:
\begin{itemize}[leftmargin=4em]
    \item At \textbf{small} and \textbf{medium} scales, the best dynamic strategies (\textbf{RL Batch} at 1262.4; \textbf{DFS Score Search} at 3824.3) \emph{exceed} the static scores (877.9 and 3003.6 respectively).
    \item At \textbf{large} scale, the static total (4407.0) appears higher in absolute terms, but it covers a different settled task count (38/39 vs.\ 32/32 for most dynamic runs), so direct comparison is limited.
    \item We suspect the gap is partly caused by \textbf{time-penalty modeling} in the MIP formulation: overdue and timing-related penalty terms may push the solver toward conservative routes that score poorly under our unified dynamic scoring function, and in some cases yield solutions worse than online heuristics.
\end{itemize}

\section{Cross-Scale Summary}

Table~\ref{tab:cross-scale-best} highlights the best-performing dynamic strategy per scale.

\begin{table}[H]
    \centering
    \caption{Best Dynamic Strategy per Scale}
    \label{tab:cross-scale-best}
    \reporttableset{0.98\textwidth}
    \begin{reporttabular}{llcc>{\raggedright\arraybackslash}l}
        \toprule
        \textbf{Scale} & \textbf{Best Algorithm} & \textbf{Score} & \textbf{On-time} & \textbf{Notes} \\
        \midrule
        Small  & RL Batch            & 1262.4 & 73.3\% & Mean of 2 seeds; DFS best on-time (81.7\%) \\
        Medium & DFS Score Search    & 3824.3 & 53.1\% & Search-based batching clearly wins \\
        Large  & Cluster Auction MAS & 3278.7 & 96.8\% & One timeout; SA Score Search best balanced \\
        \bottomrule
    \end{reporttabular}
\end{table}






Key findings across all experiments:

\begin{itemize}[leftmargin=4em]
    \item \textbf{No single algorithm wins everywhere.} Search-based loaders excel at medium scale; auction-based coordination peaks at large scale; RL/SA variants are strong at small scale.
    \item \textbf{Completion vs.\ quality.} All dynamic runs achieve near-perfect completion; on-time rate and total score are the primary discriminators, especially at medium scale where on-time drops below 60\%.
    \item \textbf{Greedy baselines remain useful.} \textbf{Nearest Task First} stays in the upper tier at small and medium scales with minimal implementation cost.
    \item \textbf{Cooperative strategies are scale-dependent.} \textbf{Relay Handoff} underperforms at all scales tested; \textbf{Cluster Auction MAS} becomes competitive only when the fleet and task pool are large enough.
    \item \textbf{Static vs.\ dynamic.} Gurobi static planning does not reliably beat online dynamic policies under the same scoring rules, suggesting that model fidelity (especially time penalties) matters as much as optimality guarantees.
\end{itemize}
